\documentclass{situationist-preprint}

\title{The Universal Availability of Meaning}
\msid{EXI-01-d}
\msfamily{EXI-01}
\msdiscipline{Meaning, Death and the Absurd}
\mstype{Consequence}
\mskeywords{meaning in life; availability; moral luck; equality; objective worth}

\begin{document}
\maketitle

\begin{abstract}
EXI-01-a holds that life is meaningful to the extent that one engages with projects of objective worth, and that meaning in this sense is available to everyone. We draw out what the availability claim implies and what it does not. It implies that meaning is not hostage to luck in outcomes, to talent or to the scale of one's projects; that no one's life is condemned to meaninglessness by circumstance alone; and that what circumstance and injustice affect is the degree of meaning, not its possibility. It does not imply that every life is meaningful, that all lives are equally meaningful, or that deprivation is harmless to meaning. We answer the objections that the view is elitist and that availability so secured is trivial.
\end{abstract}

\section{The claim}

Theorem~4 of EXI-01-a states that for every agent capable of intentional action a life of positive meaning is available. Remark~5 of EXI-01-a adds that the theorem does not settle whether any particular life is meaningful, which depends on what is done with it. EXI-17-a reaches the same availability by a different route: meaning requires active engagement with something worthwhile and not achievement (Theorem~3.3 of EXI-17-a), so that meaning is not hostage to luck in outcomes (Corollary~3.6 of EXI-17-a). EXI-01-c shows that availability does not depend on the availability of projects in the ordinary sense, since sustained engagement with a worthwhile object constitutes one. We take these results as given and ask what follows.

\section{What availability implies}

\begin{proposition}[Independence from outcome]\label{p:outcome}
Whether a life is meaningful does not depend on whether its projects succeed.
\end{proposition}

This is Corollary~3.6 of EXI-17-a, and it supports the availability claim directly. If meaning required success, it would be available only to those whom the world allows to succeed, and the world's cooperation is not in anyone's power. Because meaning requires engagement, which is within the agent's power, and worth, which belongs to the object, its availability does not wait on outcomes.

\begin{proposition}[Independence from talent and scale]\label{p:talent}
A life can be meaningful though its subject has no exceptional talent and her projects are small in scope.
\end{proposition}

Definition~1 of EXI-01-a measures meaning by engagement and by the objective worth of what one engages with. Talent bears on how much one can achieve, which is not a condition of meaning; scale bears on how much value a project serves, which can affect degree but not presence. Raising a child, tending a garden that others enjoy, or teaching a skill to a neighbor are projects of objective worth, and engagement with them confers meaning on lives that will leave no public trace.

\begin{proposition}[No condemnation by circumstance]\label{p:circ}
No agent capable of intentional action is condemned by her circumstances alone to a meaningless life.
\end{proposition}

This is Theorem~4 of EXI-01-a restated as a denial, with the strengthening of EXI-01-c, which shows that even agents without access to projects in the ordinary sense have worthwhile objects of engagement within reach.

\begin{proposition}[Degree, not possibility]\label{p:degree}
Circumstances, including injustice, affect the degree of meaning a life can attain and the ease with which it can attain it, but not whether positive meaning is available.
\end{proposition}

Definition~1 of EXI-01-a makes meaning a matter of degree, and EXI-17-a notes that achievement may add to its degree (Remark~3.5 of EXI-17-a). Time, health, education and freedom determine how much of a life can be given to engagement and how wide the range of worthwhile objects within reach is. Deprivation therefore matters to meaning, and to a theory of justice that counts it. What deprivation cannot do is remove the possibility of meaning altogether.

\section{What availability does not imply}

Three inferences are invalid. First, from availability to actuality: a life in which the agent engages with nothing worthwhile, though she could have, is not meaningful (Remark~5 of EXI-01-a). Second, from availability to equality: the degree of meaning varies with the extent of engagement and the worth of its object, and lives differ in both. Third, from availability to indifference: that meaning remains possible under oppression is no reason to tolerate oppression, which reduces the degree of meaning as it reduces other goods.

\section{Objections}

\begin{objection}[Elitism]
Objective-worth views of meaning privilege the projects valued by the educated, such as art, science and political causes, and so make meaning in practice available mainly to the privileged.
\end{objection}

\begin{reply}
The objection would hold against a view that measured worth by prestige or required projects of a particular cultural form. Definition~1 requires only that the value of a project not depend on the attitudes of those engaged in it. Caring for others, honest work that serves others' needs and the pursuit of understanding at any level meet this condition, and are open to everyone. \citet{Wolf2010} insists that the objects of meaningful engagement are not restricted to high culture; the availability claim depends on this, and the cases in Proposition~\ref{p:talent} illustrate it.
\end{reply}

\begin{objection}[Triviality]
If meaning is available to everyone, the claim that a life is meaningful is cheap and of no evaluative interest.
\end{objection}

\begin{reply}
Availability is not possession, and possession comes in degrees. That positive meaning is available to everyone is compatible with its being hard to attain in large measure, and with many lives having little of it. The claim is not that meaning is cheap but that it is not reserved: no one is shut out of it by birth, luck or the world's refusal to cooperate.
\end{reply}

\begin{objection}[The absurd]
From a sufficiently detached standpoint nothing matters, so meaning is available to no one.
\end{objection}

\begin{reply}
Remark~5 of EXI-01-a notes that Assumption~2 denies the premise of this challenge rather than refuting it \citep{Nagel1971}. EXI-04 grants that life is absurd in Nagel's sense and argues that the fitting response is irony rather than despair; the detached standpoint yields the absence of a vindication, not a demonstration that nothing is worthwhile. The availability claim inherits whatever support Assumption~2 has, and the consequences drawn here are conditional on it.
\end{reply}

\section{Conclusion}

Meaning, on the account of EXI-01-a, is open to everyone because its conditions are within everyone's reach: engagement, which is up to the agent, and worth, which belongs to what she engages with. It is not secured by success, talent, scale or fortunate circumstance, and it cannot be taken away by their absence, though it can be diminished. The availability of meaning is universal; its degree is not, and how much of it a life attains depends on what is done with it and on the conditions under which it is done \citep{Metz2013}.

\begin{thebibliography}{9}\small
\bibitem[Metz(2013)]{Metz2013} Metz, T. (2013). \emph{Meaning in Life: An Analytic Study}. Oxford University Press.
\bibitem[Nagel(1971)]{Nagel1971} Nagel, T. (1971). The absurd. \emph{The Journal of Philosophy}, 68(20), 716--727.
\bibitem[Wolf(2010)]{Wolf2010} Wolf, S. (2010). \emph{Meaning in Life and Why It Matters}. Princeton University Press.
\end{thebibliography}

\end{document}
